\section{Introduction}

First-line osimertinib has transformed care for EGFR-mutant
non-small cell lung cancer (NSCLC), yet essentially every patient
eventually relapses from residual disease
\citep{ramalingam2020,piotrowska2018}. The residual cells, known as
drug-tolerant persisters, survive targeted therapy without acquiring
resistance mutations \citep{sharma2010,hata2016}. Persisters were
initially described as a single quiescent state, but lineage-tracing
and single-cell transcriptomic studies have since shown that the
tolerant population is transcriptionally heterogeneous, spanning
several co-existing programs as cells emerge under continuous drug
pressure \citep{oren2021,aissa2021,shaffer2017}. Characterising the
geometry of this heterogeneous cell-state space---and distinguishing
a genuinely reorganised state space from one that has merely become
more variable---is central to understanding how persisters arise.

Standard single-cell trajectory tools are not designed for this
question. Branching-graph methods such as Monocle and Slingshot,
connectivity-graph methods such as PAGA, and vector-field methods
such as RNA velocity all model one-directional progression along a
tree or hierarchy \citep{saelens2019,wolf2018}. By construction they
summarise a state space as a set of connected branches and therefore
cannot quantify whether that space contains non-tree-like structure:
multi-scale features that no acyclic graph can represent. When a
minority of persister cells occupies a region of expression space
whose connectivity is not captured by any spanning tree, trajectory
inference reports a branch or a gap rather than a measurement of the
feature itself.

Persistent homology, a branch of computational topology, is built to
measure exactly such non-tree-like structure. It tracks how
$k$-dimensional features appear and disappear across a continuous
range of spatial scales and assigns each feature a persistence---the
range of scales over which it survives---yielding a scale-free,
coordinate-free summary of a point cloud's shape
\citep{edelsbrunner2010,cohen-steiner2007}. The first-homology
($H_1$) persistence quantifies the robustness of one-dimensional
structure that a connectivity tree cannot encode, and the Mapper
algorithm \citep{singh2007} provides a complementary interpretable
graph summary. Both have been applied to bulk cancer transcriptomes
\citep{nicolau2011} and to developmental single-cell data
\citep{rizvi2017}.

A critical limitation, however, has prevented persistent homology
from yielding defensible biological claims on single-cell data: prior
applications have not controlled for dispersion. The persistence of
$H_1$ features grows with the spread and anisotropy of a point cloud
even in the complete absence of organised structure, because a more
dispersed or more elongated sample of points will, by chance, enclose
larger empty regions at larger scales. Drug treatment changes the
variance and covariance of single-cell expression profiles for many
reasons unrelated to any reorganisation of the state space---altered
cell-cycle composition, shifts in library size, and broad
transcriptional adaptation. Reporting that a topological statistic
rises under treatment, or differs between conditions, is therefore
uninterpretable unless it is measured against a null that holds
dispersion fixed. Without such a control, an apparent topological
signal cannot be distinguished from a simple increase in variance,
and existing TDA-on-scRNA analyses do not make this distinction.

Here we address this gap directly. Our central methodological
contribution is a covariance-matched Gaussian null: for each cell
population we draw surrogate point clouds from a multivariate Gaussian
with the same mean and covariance as the data but no real structure,
and we ask whether the observed maximum $H_1$ persistence exceeds what
that dispersion alone produces. This converts persistent homology from
a descriptive statistic into a controlled readout that isolates
genuinely non-Gaussian topological reorganisation from changes in
spread. We pair the null with bootstrap subsampling to report
distributions rather than single point estimates, and with
cell-cycle controls.

Applying this framework, we find that in osimertinib-treated PC9
cells the maximum $H_1$ persistence is indistinguishable from the
dispersion-matched null at baseline (ratio $0.95$ at day~0, $0.91$ at
day~3) but rises above it specifically as drug-tolerant cells emerge
(ratio $1.31$ at day~7, $1.42$ at day~14), exceeding the null in
$100\%$ and $98\%$ of bootstrap subsamples respectively; the
day-7-over-day-0 and day-14-over-day-0 increases hold in $100\%$ of
bootstrap pairs. The same pattern appears in an EGFR-mutant
patient-derived xenograft, where residual disease shows a strong
excess over the null (ratio $1.95$, exceeding it in $99\%$ of
subsamples) relative to a marginal untreated baseline (ratio $1.21$),
with the residual-over-untreated increase holding in $97\%$ of
bootstrap pairs. The signal is sparse---a persistent loop occupied by
a minority of cells---and, as we show through a circular-coordinate
negative control, cells do not traverse it (Rayleigh
$R \geq 0.965$ across cohorts, versus $0.881$ for a control loop
genuinely occupied around its circumference). We therefore describe a
non-traversed topological reorganisation of the persister state space,
not population-wide cyclic motion. Erlotinib-treated cells show no
excess over the null (ratios $0.93$, $0.99$, $0.70$ at days~0, 9 and
11), which we report as specificity evidence rather than a second
positive finding, and we show that patient tumours exhibit excess
topological structure already at the treatment-naive stage (ratios
$1.52$--$1.69$ across stages), indicating that in patients this
readout reflects baseline tumour heterogeneity rather than
drug induction. All analyses are reproducible from raw data via a
single command under an open-source license, providing a
generalisable, dispersion-controlled topological readout of
single-cell state-space reorganisation.
